\section{Теоретические основы анализа производительности}

\subsection{Методология и классификация тестов производительности}

Анализ производительности вычислительных систем — это процесс сбора, измерения и интерпретации данных о работе аппаратного и программного обеспечения с целью выявления узких мест (bottlenecks) и последующей оптимизации. Центральным элементом такого анализа является \textbf{бенчмаркинг} — использование стандартизированных тестов (бенчмарков) для оценки производительности.

Бенчмарки можно классифицировать по нескольким признакам:
\begin{itemize}
    \item \textbf{Синтетические тесты (Synthetic Benchmarks):} Это специально разработанные программы для измерения пиковой производительности конкретного аппаратного компонента. Они выполняют повторяющиеся базовые операции, например, арифметические вычисления для ЦП или последовательное чтение для ОЗУ. Преимущество таких тестов в их изолированности и повторяемости. Недостаток — результаты могут не отражать реальную производительность в комплексных приложениях. В данной работе большинство тестов являются синтетическими.
    \item \textbf{Тесты-ядра (Kernel Benchmarks):} Представляют собой фрагменты кода, извлеченные из реальных приложений и выполняющие наиболее критичную к производительности функцию (например, алгоритм сжатия данных или матричное умножение).
    \item \textbf{Прикладные тесты (Application Benchmarks):} Используют реальные приложения (например, компиляторы, графические редакторы, игры) для измерения времени выполнения типичного пользовательского сценария. Это наиболее репрезентативный вид тестирования, но он сложен в автоматизации и анализе.
\end{itemize}

При проведении бенчмаркинга важно учитывать ряд факторов, способных повлиять на результат:
\begin{itemize}
    \item \textbf{Состояние системы:} Фоновые процессы, температура компонентов и текущая загрузка системы могут вносить «шум» в измерения.
    \item \textbf{Влияние кэширования:} Результаты первого запуска теста могут отличаться от последующих из-за «прогрева» кэшей процессора и дисковых кэшей ОС.
    \item \textbf{Оверхед измерения:} Сам процесс измерения времени занимает процессорные такты и может незначительно искажать результат, особенно на очень коротких операциях.
\end{itemize}
Для минимизации этих факторов в данной работе применяется подход многократного выполнения каждого теста с последующим усреднением результатов и отбрасыванием аномальных значений.

\subsection{Обзор и обоснование выбора технологического стека}

Выбор правильных инструментов является ключевым фактором успеха при разработке любого программного обеспечения. Для данного проекта основными критериями были: высокая производительность, возможность низкоуровневого взаимодействия с системой, кросс-платформенность и простота разработки графического интерфейса. На основе этих критериев был выбран стек, состоящий из языка Go и библиотеки Fyne.

\subsubsection{Язык программирования Go}

Go (Golang) — компилируемый, статически типизированный язык, разработанный в Google. Его выбор обусловлен следующими фундаментальными преимуществами:

\paragraph{Производительность и управление памятью.} Go компилируется в машинный код, обеспечивая производительность, сравнимую с C++, что является критическим требованием для ПО, измеряющего производительность. В Go реализован эффективный сборщик мусора (Garbage Collector, GC), работающий конкурентно с основной программой. Он минимизирует паузы (stop-the-world pauses), что особенно важно для отзывчивости графического интерфейса и точности измерений, не прерываемых на длительное время работой GC.

\paragraph{Модель параллелизма.} В основе параллелизма в Go лежат \textbf{горутины} и \textbf{каналы}. Горутины — это легковесные потоки, управляемые средой выполнения Go, а не ядром ОС. Их создание и переключение между ними на порядки дешевле, чем у системных потоков. Это позволяет создавать десятки тысяч горутин для эффективного распараллеливания задач. Каналы предоставляют типобезопасный механизм для коммуникации между горутинами, что позволяет избежать типичных проблем синхронизации, таких как состояние гонки (race condition). В данном проекте горутины используются для выполнения тестов в фоне, сбора данных мониторинга и обновления GUI без блокировки основного потока.

\paragraph{Кросс-компиляция и статическая линковка.} Инструментарий Go из коробки поддерживает кросс-компиляцию, позволяя с одной машины собирать исполняемые файлы для Windows, macOS, Linux и других ОС. Приложения на Go по умолчанию статически линкуются, то есть все необходимые зависимости включаются в один исполняемый файл. Это значительно упрощает распространение и запуск приложения на конечных машинах, так как не требует установки дополнительных библиотек.

\subsubsection{Библиотека для графического интерфейса Fyne}

Fyne — это кросс-платформенный GUI-инструментарий для Go, написанный на чистом Go. В сравнении с альтернативами, такими как \textbf{Gio} (ориентирован на немедленный режим рендеринга) или \textbf{Wails} (использует веб-технологии для фронтенда), Fyne был выбран за оптимальный баланс между простотой, внешним видом и производительностью.

\begin{itemize}
    \item \textbf{Декларативный подход и архитектура.} Fyne использует декларативный подход к построению интерфейса, где разработчик описывает, *что* должно быть на экране, а не *как* это нарисовать. Все виджеты являются частью дерева объектов, и Fyne автоматически управляет их рендерингом и обновлением. Рендеринг происходит с использованием аппаратного ускорения (OpenGL/Vulkan/Metal), что обеспечивает высокую производительность.
    \item \textbf{Набор виджетов и кастомизация.} Библиотека предоставляет богатый набор стандартных виджетов (кнопки, метки, прогресс-бары, карточки, вкладки), которые были активно использованы в проекте. Fyne также поддерживает темизацию, что позволяет легко изменять внешний вид приложения.
    \item \textbf{Управление жизненным циклом.} Fyne берет на себя управление главным циклом приложения и обработку событий, позволяя разработчику сосредоточиться на логике приложения.
\end{itemize}

\subsubsection{Библиотека для сбора системной информации gopsutil}

Для получения информации о системе была выбрана библиотека \texttt{gopsutil}. Это порт популярной Python-библиотеки \texttt{psutil}. Она предоставляет унифицированный кросс-платформенный API для доступа к таким метрикам, как загрузка ЦП, использование памяти, дисковая и сетевая активность, информация о процессах и многое другое. Использование \texttt{gopsutil} позволило абстрагироваться от сложностей системных вызовов в различных операционных системах и значительно ускорило разработку модуля мониторинга.

